\section{2023--2026: The Rescript, the Visitation, and a New Pontificate}

\subsection{The rescript of 20 February 2023}

On 20 February 2023 the Holy Father, in an audience granted to Cardinal
Arthur Roche, Prefect of the Dicastery for Divine Worship and the
Discipline of the Sacraments, confirmed three things about the
implementation of \work{Traditionis custodes}: that two dispensations are
``reserved in a special way to the Apostolic See (cf.\ C.I.C.\ can.\ 87
§\,1)''---the use of a parish church or the erection of a personal parish
for celebration with the 1962 Missal, and permission for priests ordained
after the motu proprio to celebrate with it; that a diocesan bishop who
has already granted such dispensations ``is obliged to inform'' the
Dicastery; and that the \work{Responsa ad dubia} of 4 December 2021 with
their explanatory notes stand confirmed. The rescript was ordered
published in \work{L'Osservatore Romano} and subsequently in \work{Acta
Apostolicae Sedis}.\endnote{Dicastery for Divine Worship and the
Discipline of the Sacraments, \latin{Rescriptum ex audientia Ss.mi} of 20
February 2023, quoted from the Holy See's own English text
(\url{https://www.vatican.va/roman_curia/congregations/ccdds/documents/rc_con_ccdds_doc_20230220_rescriptum-traditioniscustodes_en.html}),
read 2026-07-25; the exact bytes read match the artifact registered in
this repository's source library.}

The rescript does not name the Fraternity and does not purport to touch a
faculty granted by the pope to a society of apostolic life. Its effect on
the Fraternity is indirect and, on the record reached here, unquantified:
the two reserved dispensations concern precisely the two things a
flourishing traditional apostolate most often needs---a church that is a
parish church, and young priests. The Fraternity's own seminaries ordain
about fourteen men a year, and every man it has ordained since 2022 is a
priest ``ordained after the publication'' of the motu proprio; the
question whether art.~4 of \work{Traditionis custodes} reaches men
ordained for a society whose members hold a papal faculty by the decree of
11 February 2022 is not answered by any document reached for this
account.\endnote{The ordination figures are the Fraternity's own, from its
statistics page for 1 November 2025 (eleven in 2021, fifteen in 2022,
fourteen in 2023, eighteen in 2024, eleven in 2025; a twelve-year average
of fourteen). The question posed is this account's and is left open
deliberately: art.~4 addresses ``priests ordained after the publication of
the present Motu Proprio, who wish to celebrate using the \work{Missale
Romanum} of 1962,'' and requires a formal request to the diocesan bishop,
who is to consult the Apostolic See. The decree of 11 February 2022 grants
the faculty to ``each and every member,'' which on its face includes
members ordained after it. Nothing located resolves the relation between
the two.}

\subsection{The apostolic visitation opened in September 2024}

On 26 September 2024 the Fraternity published a communiqué from Fribourg
stating that it ``has recently been informed by the Dicastery for
Institutes of Consecrated Life and Societies of Apostolic Life of the
opening of an apostolic visitation of the Fraternity.'' It added, in
terms plainly meant to reassure, that ``as the Prefect of this Dicastery
himself made clear to the Superior General and his assistants during a
meeting in Rome, this visit does not originate in any problems of the
Fraternity, but is intended to enable the Dicastery to know who we are,
how we are doing and how we live, so as to provide us with any help we may
need''; that the last ordinary apostolic visit was in 2014 by the
Commission \work{Ecclesia Dei}; and that the Dicastery had been in charge
of the Fraternity and other former \work{Ecclesia Dei} institutes for the
past three years.\endnote{``Apostolic visitation of the Fraternity''
(fssp.org), Fribourg, 26 September 2024, published 27 September, read
2026-07-25. Party statement, including its account of what the Prefect
said in a private meeting.}

A statement of the Dicastery followed on 30 September 2024. This account
did not locate it in the Holy See Press Office bulletin or on a Holy See
site; it is reported by news agencies, which quote it as saying that the
Dicastery ``has called for an apostolic visitation to the Priestly
Fraternity of St.\ Peter \ldots in order to deepen the understanding of
this society of apostolic life of pontifical right and to offer the most
appropriate support to its journey of following Christ,'' and that the
visitation takes place ``in the context of the process of accompanying the
Institutes of Consecrated Life and the Societies of Apostolic Life that
were previously established by the Pontifical Commission Ecclesia Dei and
which Pope Francis' motu proprio Traditionis Custodes has placed under
the jurisdiction of this dicastery.'' The same reports state that the text
was signed by the prefect, Cardinal João Braz de Aviz, and by the
dicastery's secretary, Sister Simona Brambilla.\endnote{Agency report of
1 October 2024 carried by the \work{National Catholic Register} from
Catholic News Agency, ``Vatican announces apostolic visit to Priestly
Fraternity of St.\ Peter,'' read 2026-07-25. Bounded negative: the
Dicastery's statement of 30 September 2024 was not located in the Holy See
Press Office bulletin or elsewhere on a Holy See site, and the quotations
above are therefore quotations at one remove, from a news agency's
rendering. They are used for the stated ground of the visitation, which
matters to this account's argument, and the limitation is recorded because
it is real. The same report places the Denton seminary in Texas; it is in
Nebraska.}

If that rendering is accurate, the Dicastery gave as its ground exactly
the chain this account has traced: institutes previously established by
the Commission, placed under this dicastery by \work{Traditionis
custodes}. The visitation is thus the fourth act in the sequence begun by
faculty 6 of 1988---and, like the three before it, it is an act of the
supervisor about its own supervision.

At the cutoff of this account no report, decree or communiqué concluding
the visitation had been published by the Holy See or by the
Fraternity.\endnote{Bounded negative, checked 2026-07-25. The Fraternity's
English, French and German news feeds each carry nothing later than 16
February 2026, and their last substantive items are the audience of
January 2026 and the annual consecration novena. No item concerning the
Fraternity was located in the Holy See Press Office bulletin. A non-public
act, or an act communicated only to the institute, would not appear in
either channel.}

\subsection{The audience of 19 January 2026}

Leo XIV received the superior general, Father John Berg, in private
audience at the Vatican on Monday 19 January 2026. Father Berg was
accompanied by Father Josef Bisig, a founder, former superior general and
then rector of the seminary at Denton. The Fraternity's communiqué of 20
January describes ``the cordial half-hour meeting'' as ``an opportunity to
present to the Holy Father in greater detail the foundation and history of
the Fraternity, as well as the various forms of apostolate that it has
been offering to the faithful for almost 38 years,'' at which ``the proper
law and charism that guide the sanctification of its members were
recalled.'' The audience ``also provided an opportunity to evoke any
misunderstandings and obstacles that the Fraternity encounters in certain
places and to answer questions from the Supreme Pontiff.'' At the end the
pope gave his blessing, extending it to all members.\endnote{``Audience
with Pope Leo XIV'' (fssp.org), Fribourg, 20 January 2026, read
2026-07-25. Party account of a private audience. The photographs on the
page are credited to Vatican Media; no Holy See text about the audience
was located, and the Holy See Press Office bulletin was not found to carry
an item on it.}

The communiqué is notable for what it does not contain. It records no
papal statement about the 2022 decree, none about the visitation, and none
about \work{Traditionis custodes}. It is the most cautiously drafted of
the Fraternity's audience communiqués, and it is the only one of the
three that does not report the pope saying anything.

\subsection{July 2026: the second rupture, and the Fraternity's silence}

On 1 July 2026 the Society of Saint Pius X consecrated four bishops at
Écône without a pontifical mandate, and on 2 July the Dicastery for the
Doctrine of the Faith issued a decree declaring six named men
excommunicated and an explanatory note declaring the Society's sacred
ministers schismatic. Those acts belong to the companion volume in this
series and are not analysed here.\endnote{Dicastery for the Doctrine of
the Faith, \latin{Decreto} and \latin{Nota Esplicativa}, Prot.\
N.~99/2009, 2 July 2026, Holy See Press Office bulletin B0568
(\url{https://press.vatican.va/content/salastampa/it/bollettino/pubblico/2026/07/02/0568/01077.html}
and \ldots/01078.html), read 2026-07-25; the exact bytes read match the
artifacts registered in this repository's source library. The events and
their canonical analysis are the subject of TPI-01 in this series.}

They are recorded here for one reason. The Fraternity of Saint Peter
exists because, in 1988, a group of priests judged that consecrating
bishops without a papal mandate was not a course they could follow, and
said so in writing on the second day after it happened. In 2026 the same
act was performed again, by the same body, and the Fraternity's general
house published nothing. Its English, French and German news feeds are
silent from 16 February 2026 to the cutoff of this
account.\endnote{Bounded negative, checked 2026-07-25 in the Fraternity's
three language feeds at fssp.org. Silence in a published feed is not
silence \latin{tout court}: the institute may have communicated internally,
through its provinces and districts, or to the Holy See. What can be said
is that the general house issued no public communiqué in its ordinary
channel. This account draws no inference about the reason.}

\begin{documentinventory}
\inventoryrow{Rescript from the audience of 20 February 2023}{Holy See
English text}{The two dispensations reserved to the Apostolic See, the
bishops' duty to inform, and confirmation of the 2021 \work{Responsa}. The
Fraternity is not named.}
\inventoryrow{Fraternity communiqué, 26 September 2024}{fssp.org}{That an
apostolic visitation was opened by the dicastery for consecrated life, and
the Fraternity's account of the Prefect's assurances.}
\inventoryrow{Dicastery statement, 30 September 2024}{Not located; quoted
by a news agency}{The stated ground of the visitation, at one remove.}
\inventoryrow{Outcome of the visitation}{Not published; not
located}{Nothing. No report, decree or communiqué concluding it was found
through 25 July 2026.}
\inventoryrow{Fraternity communiqué, 20 January 2026}{fssp.org}{That Leo
XIV received the superior general and a founder in private audience on 19
January 2026 and gave his blessing. No papal statement is reported.}
\inventoryrow{Fraternity news feeds, 16 February to 25 July
2026}{fssp.org, three languages}{Bounded negative: no public communiqué
of any kind, including on the consecrations of 1 July 2026.}
\end{documentinventory}
